%! Independence, Basis, and Dimension — Unit 3, Lesson 3.4 — Student Packet Cover Sheet
\documentclass[10pt]{article}
\usepackage{linalg-article}
\usepackage{linalg-boxes}
\usepackage{ltablex}
\keepXColumns

\begin{document}

% Full-bleed burgundy banner
\begin{tikzpicture}[remember picture, overlay]
  \fill[burgundy] (current page.north west)
    rectangle ([yshift=-0.9in]current page.north east);
\end{tikzpicture}
\vspace{-0.8in}

\begin{center}
  {\color{white}\textbf{\LARGE Linear Algebra}}\\[4pt]
  {\color{blushmid}\large\textbf{Unit 3 \quad The Four Fundamental Subspaces}}\\[6pt]
  {\color{white}\textbf{\large Lesson 3.4 \quad Independence, Basis, and Dimension}}
\end{center}

\vspace{0.22in}
\namedateperiod
\vspace{0.18in}

\begin{learningtargetbox}
\small
\begin{enumerate}[leftmargin=1.5em, itemsep=5pt]
  \item \ldots{} test a set of vectors for \textbf{linear independence} by checking whether $A\mathbf{x}=\mathbf{0}$ has only the trivial solution ($N(A)=\{\mathbf{0}\}$).
  \item \ldots{} build a \textbf{basis} --- an independent set that spans a space: the pivot columns for the column space $C(A)$ and the special solutions for the nullspace $N(A)$.
  \item \ldots{} give the \textbf{dimension} of a space ($\dim C(A)=r$, $\dim N(A)=n-r$) and explain what that number counts.
\end{enumerate}
\end{learningtargetbox}

\vspace{0.14in}

\begin{tocbox}
\small
\renewcommand{\arraystretch}{1.5}
\begin{tabularx}{\linewidth}{@{} c l X r @{}}
\rowcolor{blush}
\textbf{\#} & \textbf{Component} & \textbf{Description} & \textbf{Score} \\
\midrule
1 & Warm-Up        & Spiral review: parallel vectors, a nullspace special solution, pivot vs.\ free columns & \blank{1.2cm} \\
2 & Guided Notes   & Independence via $N(A)$, basis (independent $+$ spans), dimension $=r$ and $n-r$     & \blank{1.2cm} \\
3 & Group Activity & Independent, Basis, Dimension --- reading all three from one reduction              & \blank{1.2cm} \\
4 & Exit Ticket    & Quick check on independence, a basis with its dimension, and what dimension counts   & \blank{1.2cm} \\
5 & Homework       & Independence tests, bases \& dimensions of $C(A)$ and $N(A)$, four-subspaces preview  & \blank{1.2cm} \\
\midrule
  & \multicolumn{2}{r}{\textbf{Total}} & \blank{1.2cm} \\
\end{tabularx}
\end{tocbox}

\vspace{0.10in}

\begin{remindbox}
\small A set of vectors is \textbf{linearly independent} when the only combination equal to $\mathbf{0}$ is
the all-zero one --- equivalently, stacking them as the columns of $A$, when $A\mathbf{x}=\mathbf{0}$ forces
$\mathbf{x}=\mathbf{0}$ ($N(A)=\{\mathbf{0}\}$). A \textbf{basis} is an independent set that also
\emph{spans} the space --- just enough vectors to reach everything, with none wasted. The pivot columns of
$A$ are a basis for $C(A)$; the special solutions are a basis for $N(A)$. The \textbf{dimension} is the
number of vectors in a basis (the same for every basis): $\dim C(A)=r$ and $\dim N(A)=n-r$. Dimension
counts the independent directions --- the ``degrees of freedom'' --- in the space.
\end{remindbox}

\end{document}
